% GENERATED FILE. Edit canonical lessons, then run npm run generate:books.
% canonical-source-hash: fnv1a64:06c6d8d5cf91c028
% canonical-lessons: SA-C338-punyam, SA-C338-papam, SA-C338-svargah, SA-C338-narakah, SA-C338-bhunkte

\chapter{Heaven and Hell}
\label{ch:sa-heaven-and-hell}
\hlchaptermodality{338}

\begin{chapteropening}
\textbf{By the end of this chapter:} \emph{I can say what the old stories teach about good and bad deeds.}

Complete the last lesson of chapter 338: I can say what the old stories teach about good and bad deeds.
\end{chapteropening}

\section[\texorpdfstring{\sk{पुण्यम्} (puṇyam)}{पुण्यम् (puṇyam)}]{\sk{पुण्यम्} (puṇyam) — merit; a good deed}

\label{lesson:SA-C338-punyam}



\begin{quote}
Before the new one: say the Sanskrit for promises, then the Sanskrit for hunting.
\end{quote}

\subsection*{You'll want to know: \sk{पुण्यम्}}
\textbf{\sk{पुण्यम्}} — \emph{puṇyam} — \textquotedblleft{}merit; a good deed\textquotedblright{}.

\begin{quote}
\sk{दानेन} \sk{पुण्यं} \sk{वर्धते।}
\end{quote}

\begin{quote}
\emph{dānena puṇyaṃ vardhate.}
\end{quote}

\begin{quote}
Merit grows through giving.
\end{quote}

A neuter noun.

\begin{cousinweb}
\textbf{\sk{पुण्यम्}} (\emph{puṇyam}) MW suggests, with a perhaps, the root √puṣ, \textquotedblleft{}to thrive\textquotedblright{}, and gives \emph{puṇya} as good, virtuous and, as a noun, virtue, a good work, religious merit, so its origin is uncertain (MW).
\end{cousinweb}

\begin{grammarlens}[title={What you've built}]
One more word for telling what happened.
\end{grammarlens}

\subsection*{Your turn}
\begin{itemize}
\raggedright
  \item \emph{Say it:} \emph{puṇyam}
  \item \emph{Say it:} \emph{puṇyam}, once more
  \item \emph{Say it:} \emph{dānena puṇyaṃ vardhate.} — Merit grows through giving.
  \item \emph{Recall:} say the Sanskrit for promises, then the Sanskrit for hunting, then say \emph{puṇyam} again
\end{itemize}

\begin{tcolorbox}[breakable,colback=teal!4,colframe=teal!35!black,title={Before you move on}]
What does \sk{पुण्यम्} mean? (\textquotedblleft{}Merit; a good deed\textquotedblright{}.) Say it once more.
\end{tcolorbox}
\section[\texorpdfstring{\sk{पापम्} (pāpam)}{पापम् (pāpam)}]{\sk{पापम्} (pāpam) — a sin; evil}

\label{lesson:SA-C338-papam}



\begin{quote}
Before the new one: say the Sanskrit for hunting, then the Sanskrit for merit.
\end{quote}

\subsection*{You'll want to know: \sk{पापम्}}
\textbf{\sk{पापम्}} — \emph{pāpam} — \textquotedblleft{}a sin; evil\textquotedblright{}.

\begin{quote}
\sk{चौर्यं} \sk{पापम्} \sk{अस्ति।}
\end{quote}

\begin{quote}
\emph{cauryaṃ pāpam asti.}
\end{quote}

\begin{quote}
Theft is a sin.
\end{quote}

A neuter noun.

\begin{cousinweb}
\textbf{\sk{पापम्}} (\emph{pāpam}) MW gives \emph{pāpa} as bad, wicked, evil and, as a noun, evil, sin, without a root in the entry, so its origin is left uncertain here (MW).
\end{cousinweb}

\begin{grammarlens}[title={What you've built}]
One more word for telling what happened.
\end{grammarlens}

\subsection*{Your turn}
\begin{itemize}
\raggedright
  \item \emph{Say it:} \emph{pāpam}
  \item \emph{Say it:} \emph{pāpam}, once more
  \item \emph{Say it:} \emph{cauryaṃ pāpam asti.} — Theft is a sin.
  \item \emph{Recall:} say the Sanskrit for hunting, then the Sanskrit for merit, then say \emph{pāpam} again
\end{itemize}

\begin{tcolorbox}[breakable,colback=teal!4,colframe=teal!35!black,title={Before you move on}]
What does \sk{पापम्} mean? (\textquotedblleft{}A sin; evil\textquotedblright{}.) Say it once more.
\end{tcolorbox}
\section[\texorpdfstring{\sk{स्वर्गः} (svargaḥ)}{स्वर्गः (svargaḥ)}]{\sk{स्वर्गः} (svargaḥ) — heaven}

\label{lesson:SA-C338-svargah}



\begin{quote}
Before the new one: say the Sanskrit for merit, then the Sanskrit for a sin.
\end{quote}

\subsection*{You'll want to know: \sk{स्वर्गः}}
\textbf{\sk{स्वर्गः}} — \emph{svargaḥ} — \textquotedblleft{}heaven\textquotedblright{}.

\begin{quote}
\sk{पुण्येन} \sk{जनः} \sk{स्वर्गं} \sk{गच्छति।}
\end{quote}

\begin{quote}
\emph{puṇyena janaḥ svargaṃ gacchati.}
\end{quote}

\begin{quote}
Through merit a person goes to heaven.
\end{quote}

A masculine noun.

\begin{cousinweb}
\textbf{\sk{स्वर्गः}} (\emph{svargaḥ}) is built from \emph{svar}, \textquotedblleft{}light, the sky\textquotedblright{}, and \emph{ga}, \textquotedblleft{}going\textquotedblright{}, from the root √gam; MW gives \emph{svarga} as leading to heaven and, as a noun, heaven (MW).
\end{cousinweb}

\begin{grammarlens}[title={What you've built}]
One more word for telling what happened.
\end{grammarlens}

\subsection*{Your turn}
\begin{itemize}
\raggedright
  \item \emph{Say it:} \emph{svargaḥ}
  \item \emph{Say it:} \emph{svargaḥ}, once more
  \item \emph{Say it:} \emph{puṇyena janaḥ svargaṃ gacchati.} — Through merit a person goes to heaven.
  \item \emph{Recall:} say the Sanskrit for merit, then the Sanskrit for a sin, then say \emph{svargaḥ} again
\end{itemize}

\begin{tcolorbox}[breakable,colback=teal!4,colframe=teal!35!black,title={Before you move on}]
What does \sk{स्वर्गः} mean? (\textquotedblleft{}Heaven\textquotedblright{}.) Say it once more.
\end{tcolorbox}
\section[\texorpdfstring{\sk{नरकः} (narakaḥ)}{नरकः (narakaḥ)}]{\sk{नरकः} (narakaḥ) — hell}

\label{lesson:SA-C338-narakah}



\begin{quote}
Before the new one: say the Sanskrit for a sin, then the Sanskrit for heaven.
\end{quote}

\subsection*{You'll want to know: \sk{नरकः}}
\textbf{\sk{नरकः}} — \emph{narakaḥ} — \textquotedblleft{}hell\textquotedblright{}.

\begin{quote}
\sk{पापेन} \sk{जनः} \sk{नरकं} \sk{गच्छति} \sk{इति} \sk{वदन्ति।}
\end{quote}

\begin{quote}
\emph{pāpena janaḥ narakaṃ gacchati iti vadanti.}
\end{quote}

\begin{quote}
People say that sin leads a person to hell.
\end{quote}

A masculine noun.

\begin{cousinweb}
\textbf{\sk{नरकः}} (\emph{narakaḥ}) MW gives \emph{naraka} as hell, a place of torment, without a derivation, so its origin is uncertain (MW).
\end{cousinweb}

\begin{grammarlens}[title={What you've built}]
One more word for telling what happened.
\end{grammarlens}

\subsection*{Your turn}
\begin{itemize}
\raggedright
  \item \emph{Say it:} \emph{narakaḥ}
  \item \emph{Say it:} \emph{narakaḥ}, once more
  \item \emph{Say it:} \emph{pāpena janaḥ narakaṃ gacchati iti vadanti.} — People say that sin leads a person to hell.
  \item \emph{Recall:} say the Sanskrit for a sin, then the Sanskrit for heaven, then say \emph{narakaḥ} again
\end{itemize}

\begin{tcolorbox}[breakable,colback=teal!4,colframe=teal!35!black,title={Before you move on}]
What does \sk{नरकः} mean? (\textquotedblleft{}Hell\textquotedblright{}.) Say it once more.
\end{tcolorbox}
\section[\texorpdfstring{\sk{भुंक्ते} (bhuṅkte)}{भुंक्ते (bhuṅkte)}]{\sk{भुंक्ते} (bhuṅkte) — reaps; enjoys}

\label{lesson:SA-C338-bhunkte}



\begin{quote}
Before the new one: say the Sanskrit for heaven, then the Sanskrit for hell.
\end{quote}

\subsection*{You'll want to know: \sk{भुंक्ते}}
\textbf{\sk{भुंक्ते}} — \emph{bhuṅkte} — \textquotedblleft{}reaps; enjoys\textquotedblright{}.

\begin{quote}
\sk{जनः} \sk{पुण्यस्य} \sk{फलं} \sk{भुंक्ते।}
\end{quote}

\begin{quote}
\emph{janaḥ puṇyasya phalaṃ bhuṅkte.}
\end{quote}

\begin{quote}
A person reaps the fruit of merit.
\end{quote}

Past (he …): \sk{सः} \sk{फलं} \sk{भुक्तवान्}, \emph{saḥ phalaṃ bhuktavān}.

\begin{cousinweb}
\textbf{\sk{भुंक्ते}} (\emph{bhuṅkte}) comes from the root √bhuj, \textquotedblleft{}to enjoy, use, eat\textquotedblright{}, the root of \emph{bhojana}, \textquotedblleft{}a meal\textquotedblright{}; MW gives it as to enjoy, use, possess, consume (MW).
\end{cousinweb}

\begin{grammarlens}[title={What you've built}]
One more word for telling what happened.
\end{grammarlens}

\subsection*{Your turn}
\begin{itemize}
\raggedright
  \item \emph{Say it:} \emph{bhuṅkte}
  \item \emph{Say it:} \emph{bhuṅkte}, once more
  \item \emph{Say it:} \emph{janaḥ puṇyasya phalaṃ bhuṅkte.} — A person reaps the fruit of merit.
  \item \emph{Recall:} say the Sanskrit for heaven, then the Sanskrit for hell, then say \emph{bhuṅkte} again
\end{itemize}

\begin{tcolorbox}[breakable,colback=teal!4,colframe=teal!35!black,title={Before you move on}]
What does \sk{भुंक्ते} mean? (\textquotedblleft{}Reaps; enjoys\textquotedblright{}.) Say it once more.
\end{tcolorbox}
